\chapter{Introducción}

En el presente informe se aborda el estudio del índice horario en
transformadores trifásicos. El índice horario es un parámetro fundamental que
define el desfase angular entre las tensiones del bobinado de alta tensión (AT)
y el de baja tensión (BT).

\section{Objetivos}
\begin{itemize}
    \item Definir el concepto de índice horario y su representación.
    \item Identificar las diferentes conexiones de transformadores y sus 
    respectivos índices.
    \item Analizar las ventajas y aplicaciones prácticas de la selección de un 
    determinado índice horario.
    \item Comprender las condiciones necesarias para la puesta en paralelo de 
    transformadores.
\end{itemize}

El correcto entendimiento de este concepto es crucial para garantizar la 
operación segura y eficiente de los sistemas de potencia, especialmente cuando 
se requiere la interconexión de diferentes máquinas o redes.
